\begin{corollary}
Fix \(d=3\) or \(4\). For each attainable state after the first pair,
choose one of the \(d\)-pair unmarked returns constructed in the proof
(for length four, a two-pair return may be repeated). At \(k\)
prescribed disjoint intervals of \(d\) pairs, all after the first
pair, the probability of seeing the chosen word at every interval
under the \((u,v)\)-weighted law is exactly
\[
                    \frac{Z_{n-dk}(u,v)}{Z_n(u,v)}
\]
whenever the placement is possible and the denominator is positive.
For fixed \(k\) this tends to \(\rho(u,v)^{dk}\).
\end{corollary}
\begin{proof}
Delete the indicated words. Each restores its full incoming state,
including the previous letter, and contributes weight one.
Conversely the prescribed positions and the state there determine
each inserted word uniquely. This bijection preserves the weight of
the rest of the diagram. The ratio limit gives the assertion. Applying
the same formula to each subset shows joint convergence of the indicators
on disjoint intervals to independent Bernoulli variables of parameter
\(\rho(u,v)^d\).
\end{proof}
